% Attack results (C5), drafted 2026-09-05 from results/memorizing_model_recall.csv, results/composition_summary.csv,
% results/certificate_caps_memoriser.csv. Prose revised 2026-09-05.
\subsection{What the attacks recover}\label{sec:attack-results}

\paragraph{The memorising model}
Without the anchor, the fine-tuned model reproduces the training passages from their first $20$ tokens with mean near-verbatim recall $0.91$ (attack split) and $0.92$ (validation) under greedy decoding, $86$ to $88\%$ of passages above $0.8$, and $0.80$ to $0.81$ under the temperature-1 sampling the anchored decoder requires; the held-out test passages are not reproduced at all (recall $0.00$, longest span two words), so the memorisation is specific to the trained text. The model assigns the attack-split passages $0.19$ nats per token where the anchor assigns $3.28$, so reproducing a passage verbatim requires the served distribution to diverge from the anchor by about $3.1$ nats per token, the same number as the budget rate $k=3$. The prefix debt charged on a memorised prompt is correspondingly large: $\delta_{\mathrm{init}} = 10.4$ nats for the single query and $12.9$ for a window with the true prefix, against $6.2$ for the non-memorising model on the CopyBench prompts of the released runs.

\paragraph{Composition}
Figure~\ref{fig:composition} reports the attack on 100 memorised passages (targets of $260$ tokens on average; $27{,}818$ budgeted queries, none over budget). Three regimes appear. For $k \le 1$ no strategy recovers anything: the budget binds (single-query utilisation $0.81$ to $0.97$), and for $20$-token windows the budget $kL$ is $3$ to $20$ nats against a prefix debt of $12.9$, so the anchor writes every token. At $k=3$, where the certificate is vacuous for every one of these passages, the decoder still reproduces only $1.2\%$ of a passage in a single query and $9.7\%$ through oracle windows of $50$ tokens, at utilisation above $0.9$: the protection comes from a budget that binds, and the certificate does not describe it. Between $k=5$ and $k=20$ the protection disappears. Single-query recall rises from $0.09$ to $0.48$, against $0.49$ for the unconstrained model, while utilisation falls to $0.18$, and oracle windows recover $0.86$ of a passage at $k=20$ against $0.88$ unconstrained. In between, composition is the better strategy: $50$-token windows with the true prefix double the single query's recall at $k=5$ ($0.23$ against $0.09$) and $k=10$ ($0.41$ against $0.20$), every query compliant. Chained windows, which use only the adversary's own outputs as prefixes, track the single query ($0.11$ at $k=10$, $0.50$ at $k=20$), because an early error derails the later windows; an adversary with the true prefix, which is the position of a rights-holder testing their own work or of an attacker who already holds most of it, does much better. He et al.\ sweep $k$ up to $20$; the audited budgets $k \in \{3, 5\}$ are the last on that axis at which the mechanism binds against a memoriser, and the certificate is uninformative at all of them.

\paragraph{Where the budget stops binding}
Proposition~\ref{prop:path} gives a condition, computable from the anchor and the target alone, that verbatim reproduction at rate $k$ must satisfy: the anchor's surprisal prefix sums of the passage, offset by the prefix debt, must never cross the budget line $(t+1)k$. For these 100 passages the critical rate has median $13.9$ nats per token, $4.3$ times the mean per-token surprisal of $3.24$, because the budget must cover the running maximum of a dispersed surplus. No passage satisfies the condition at $k \le 5$, $2\%$ do at $k=10$ and $92\%$ at $k=20$, which brackets the observed transition. Below the critical rate the bank pays for almost every token --- $87\%$ at $k=3$ and $99\%$ at $k=5$ --- while the anchor writes the first $4.8$ and $2.8$ tokens as the debt is repaid, $1.1$ at $k=10$ and none at $k=20$. What changes across the transition is therefore not the budget's ability to pay but the anchor-written opening, which derails the memorised continuation; Section~\ref{sec:prefixdebt} switches the debt off to test that.
